% Paper 3: RML - Abstract (LaTeX-ready)
% Meta-Policy Learning Over QA Manifolds
% Generated: 2025-12-29

\begin{abstract}
Learning and control on discrete state spaces governed by algebraic invariants presents a fundamental challenge: how can agents navigate complex manifolds efficiently when ground-truth queries are expensive? We introduce \textbf{structure-aware control}, a paradigm where learned topological predicates guide action selection with minimal oracle usage. Rather than learning forward dynamics or value functions, structure-aware policies learn which worlds are possible—structural properties that determine reachability. We evaluate this approach on a diagonal reachability task over the Caps(30,30) manifold using a pre-trained World Model (QAWM) that predicts return-in-$k$ from sparse samples. Our central finding is that learned queries dominate ground-truth simulation in oracle-limited regimes: QAWM-Greedy achieves \textbf{4.20 successes per oracle call} versus Oracle-Greedy's 2.97—a \textbf{1.41$\times$ efficiency advantage per success}—despite lower absolute success rates (32\% vs 60\%). An ablation study reveals that global topological structure (return-in-$k$ predictions) provides stronger control signals than local feasibility constraints (legality predictions), establishing a \textbf{topology-over-constraints principle} for planning on discrete manifolds. This work demonstrates that offline-learned structural knowledge can enable oracle-efficient control without task-specific reward engineering, with applications to combinatorial optimization, neural architecture search, and other domains where discrete dynamics exhibit learnable topological structure.
\end{abstract}
